\documentclass[11pt]{article}

\usepackage[a4paper,margin=1in]{geometry}
\usepackage[T1]{fontenc}
\usepackage{lmodern}
\usepackage{microtype}
\usepackage{amsmath,amssymb,amsthm,mathtools}
\usepackage{booktabs,longtable,array}
\usepackage{enumitem}
\usepackage{xcolor}
\usepackage[colorlinks=true,linkcolor=blue!55!black,urlcolor=blue!55!black]{hyperref}

\newtheorem{theorem}{Theorem}[section]
\newtheorem{lemma}[theorem]{Lemma}
\newtheorem{proposition}[theorem]{Proposition}
\newtheorem{corollary}[theorem]{Corollary}
\theoremstyle{definition}
\newtheorem{definition}[theorem]{Definition}
\newtheorem{remark}[theorem]{Remark}

\newcommand{\CC}{\mathbb C}
\newcommand{\RR}{\mathbb R}
\newcommand{\NN}{\mathbb N}
\newcommand{\Sn}{\mathfrak S}
\newcommand{\per}{\operatorname{per}}
\newcommand{\tr}{\operatorname{tr}}
\newcommand{\diag}{\operatorname{diag}}
\newcommand{\Rea}{\operatorname{Re}}
\newcommand{\Ima}{\operatorname{Im}}
\newcommand{\adj}{\operatorname{adj}}
\newcommand{\rank}{\operatorname{rank}}
\newcommand{\ip}[2]{\left\langle #1,#2\right\rangle}
\newcommand{\norm}[1]{\left\lVert #1\right\rVert}
\newcommand{\abs}[1]{\left\lvert #1\right\rvert}
\newcommand{\mat}[1]{\begin{pmatrix}#1\end{pmatrix}}
\newcommand{\psd}{\succeq 0}
\newcommand{\pd}{\succ 0}
\newcommand{\eps}{\varepsilon}

\title{Permanental Dominance in Dimension Four\\
  \large A Direct Geometric Proof}
\author{Siwei Zeng}
\date{August 2026}
\hypersetup{
  pdftitle={Permanental Dominance in Dimension Four: A Direct Geometric Proof},
  pdfauthor={Siwei Zeng},
  pdfsubject={The complete normalized permanental-dominance theorem in dimension four}
}

\begin{document}
\maketitle

\begin{abstract}
For a subgroup $H\leq \Sn_4$ and an irreducible complex character $\chi$ of
$H$, let
\[
  d_\chi^H(A)=\sum_{\sigma\in H}\chi(\sigma)
  \prod_{i=1}^4 A_{i,\sigma(i)}
\]
be the associated generalized matrix function.  We prove that every
Hermitian positive-semidefinite matrix $A\in M_4(\CC)$ satisfies
\[
  \frac{1}{\chi(1)}d_\chi^H(A)\leq \per A.
\]
Both sides are real.  The proof first reduces to correlation matrices and
then treats the thirty-seven irreducible-character rows belonging to the
eleven conjugacy types of subgroups of $\Sn_4$.  All but two rows follow from
elementary principal-minor, moment, and block-contraction inequalities.  The
two remaining rows are the non-real one-dimensional characters of $A_4$.
For them the difference from the permanent becomes an affine Hermitian
quadratic form.  A residual-vector normalization turns this form into a
polynomial on the cone of $3\times3$ Gram matrices.  We prove its
nonnegativity by a rank-one sum of squares, a positive-definite
two-dimensional Hessian, and exact rational Gram certificates.  The argument
also yields a positive-semidefinite adjugate certificate for the exceptional
$A_4$ rows.
\end{abstract}

\tableofcontents

\section{Statement of the theorem}

Let $n\geq1$, let $H\leq\Sn_n$, and let $\chi:H\to\CC$ be a character.
For $A=(A_{ij})\in M_n(\CC)$ put
\begin{equation}\label{eq:def-dchi}
 d_\chi^H(A):=\sum_{\sigma\in H}\chi(\sigma)
   \prod_{i=1}^n A_{i,\sigma(i)},
 \qquad
 \per A:=\sum_{\sigma\in\Sn_n}\prod_{i=1}^n A_{i,\sigma(i)}.
\end{equation}
The normalization in the permanental-dominance conjecture is
$d_\chi^H/\chi(1)$.  Notice that no complex conjugate is placed on
$\chi(\sigma)$ in~\eqref{eq:def-dchi}; this convention agrees with all
formulae below.

\begin{lemma}[Reality]\label{lem:reality}
If $A=A^*$ and $\chi$ is the character of a finite-dimensional complex
representation of $H$, then $d_\chi^H(A)$ and $\per A$ are real.
\end{lemma}

\begin{proof}
Every finite-group representation may be made unitary, and hence
$\overline{\chi(\sigma)}=\chi(\sigma^{-1})$.  Hermitian symmetry gives
\[
 \overline{\prod_i A_{i,\sigma(i)}}
 =\prod_i A_{\sigma(i),i}
 =\prod_j A_{j,\sigma^{-1}(j)}.
\]
Conjugating the sum and replacing $\sigma$ by $\sigma^{-1}$ proves the first
claim.  The second is the special case $H=\Sn_n$, $\chi=1$.
\end{proof}

\begin{theorem}[Permanental dominance for $n=4$]\label{thm:main}
Let $H\leq\Sn_4$, let $V$ be an irreducible finite-dimensional complex
representation of $H$, and let $\chi$ be its character.  For every
$A\in M_4(\CC)$ with $A\psd$,
\begin{equation}\label{eq:main}
  \frac{d_\chi^H(A)}{\dim V}\leq \per A.
\end{equation}
\end{theorem}

The scope of Theorem~\ref{thm:main} is exactly dimension four.  It includes
every subgroup of $\Sn_4$ and every irreducible complex representation of
that subgroup, but it does not assert the conjecture in arbitrary dimension.

\section{Reduction to correlation matrices}

The normalization step is elementary but important because it removes all
diagonal parameters without an invertibility hypothesis.

\begin{lemma}[Diagonal congruence]\label{lem:scaling}
Let $s_1,\dots,s_n\in\RR$ and $D=\diag(s_1,\dots,s_n)$.  Then
\[
 d_\chi^H(DAD)=\left(\prod_i s_i\right)^2d_\chi^H(A),
 \qquad
 \per(DAD)=\left(\prod_i s_i\right)^2\per A.
\]
\end{lemma}

\begin{proof}
For every permutation $\sigma$,
\[
 \prod_i(DAD)_{i,\sigma(i)}
 =\prod_i(s_iA_{i,\sigma(i)}s_{\sigma(i)})
 =\left(\prod_i s_i\right)^2\prod_iA_{i,\sigma(i)}.
\]
Summation proves both identities.
\end{proof}

\begin{lemma}[Correlation reduction]\label{lem:correlation-reduction}
It is enough to prove Theorem~\ref{thm:main} for positive-semidefinite
matrices whose diagonal entries are all one.
\end{lemma}

\begin{proof}
Suppose first that $A_{ii}>0$ for all $i$ and put
$s_i=\sqrt{A_{ii}}$.  Since a Hermitian matrix has real diagonal, the matrix
\[
 C=\diag(s_i^{-1})A\diag(s_i^{-1})
\]
is positive semidefinite and has diagonal one.  Moreover
$A=\diag(s_i)C\diag(s_i)$, so Lemma~\ref{lem:scaling} transfers the desired
inequality from $C$ to $A$.

If $A_{ii}=0$ for some $i$, positivity of the $2\times2$ principal minors
implies $A_{ij}=A_{ji}=0$ for every $j$.  Each permutation monomial then
contains a zero entry in row $i$, so both sides of~\eqref{eq:main} vanish.
\end{proof}

We therefore write a general $4\times4$ correlation matrix as
\begin{equation}\label{eq:corr-six}
 A(a,b,c,d,e,f)=
 \mat{
 1&a&b&c\\
 \bar a&1&d&e\\
 \bar b&\bar d&1&f\\
 \bar c&\bar e&\bar f&1
 }\psd.
\end{equation}

\section{The finite group-theoretic reduction}

\subsection{Subgroups and characters}

Up to conjugacy, the subgroups of $\Sn_4$ have the following eleven types.
Here $r=(1234)$ and $s=(13)$ in the dihedral row.
\begin{center}
\begin{tabular}{@{}lllrr@{}}
\toprule
type & representative & abstract group & order & irreducible rows\\
\midrule
$1$ & $\{1\}$ & $1$ & 1&1\\
$C_2^{\rm t}$ & $\langle(12)\rangle$ & $C_2$ &2&2\\
$C_2^{\rm d}$ & $\langle(12)(34)\rangle$ & $C_2$ &2&2\\
$C_3$ & $\langle(123)\rangle$ & $C_3$ &3&3\\
$C_4$ & $\langle(1234)\rangle$ & $C_4$ &4&4\\
$V_4^{\rm n}$ & $\{1,(12)(34),(13)(24),(14)(23)\}$ & $C_2^2$ &4&4\\
$V_4^{\rm d}$ & $\langle(12),(34)\rangle$ & $C_2^2$ &4&4\\
$S_3$ & $\operatorname{Stab}(4)$ & $S_3$ &6&3\\
$D_8$ & $\langle r,s\rangle$ & $D_8$ &8&5\\
$A_4$ & $A_4$ & $A_4$ &12&4\\
$S_4$ & $S_4$ & $S_4$ &24&5\\
\bottomrule
\end{tabular}
\end{center}
The row counts sum to $37$.

\begin{lemma}[Subgroup classification]\label{lem:subgroup-classification}
Every subgroup of $\Sn_4$ is conjugate to exactly one of the eleven types in
the preceding table.
\end{lemma}

\begin{proof}
We recall the elementary classification.  An intransitive subgroup preserves
an orbit partition of $4$.  The partitions $1+1+1+1$, $2+1+1$, $2+2$ and
$3+1$ give respectively $1$, the two possible embedded copies of $C_2$,
$V_4^{\rm d}$ and subgroups of the displayed $S_3$.  This yields
$1,C_2^{\rm t},C_2^{\rm d},C_3,V_4^{\rm d},S_3$.

Now suppose the action is transitive.  Orbit--stabilizer shows that the order
is divisible by $4$ and divides $24$.  A Sylow analysis gives the following
possibilities.  A cyclic Sylow $2$-subgroup generated by a four-cycle gives
$C_4$; a non-cyclic subgroup of order $4$ is the normal double-transposition
group $V_4^{\rm n}$; a subgroup of order $8$ is a Sylow $2$-subgroup and is
conjugate to $\langle(1234),(13)\rangle$; order $12$ gives the unique
index-two subgroup $A_4$; and order $24$ gives $S_4$.  These groups have
different order, orbit structure, or involution type, so the listed types are
not conjugate to one another.
\end{proof}

For later reference, the non-cyclic character tables are recorded below.
The class order for $S_4$ is
$1,(12),(12)(34),(123),(1234)$:
\begin{equation}\label{eq:s4-table}
\begin{array}{c|rrrrr}
 &1&(12)&(12)(34)&(123)&(1234)\\\hline
{[4]}    &1& 1& 1& 1& 1\\
{[1^4]}  &1&-1& 1& 1&-1\\
{[31]}   &3& 1&-1& 0&-1\\
{[211]}  &3&-1&-1& 0& 1\\
{[22]}   &2& 0& 2&-1& 0
\end{array}
\end{equation}
For $A_4$, with $C_+$ and $C_-$ the two classes of three-cycles and
\[
 \omega=-\frac12-\frac{\sqrt3}{2}i,
 \qquad \omega^2+\omega+1=0,
\]
the table is
\begin{equation}\label{eq:a4-table}
\begin{array}{c|rrrr}
 &1&(12)(34)&C_+&C_-\\\hline
1&1&1&1&1\\
\chi_\omega&1&1&\omega&\omega^2\\
\bar\chi_\omega&1&1&\omega^2&\omega\\
\chi_3&3&-1&0&0.
\end{array}
\end{equation}
For $D_8=\langle r,s:r^4=s^2=1,srs=r^{-1}\rangle$, the four linear
characters take independent signs on $r$ and $s$, while its two-dimensional
character is $2,-2,0,0,0$ on
$1,r^2,\{r,r^3\},\{s,r^2s\},\{rs,r^3s\}$.  The familiar tables of $C_m$,
$C_2^2$, and $S_3$ complete all thirty-seven rows.  Orthogonality of the
displayed rows and the identity $\sum_\chi\chi(1)^2=|H|$ prove that no
irreducible character is missing.

\subsection{Conjugation invariance}

\begin{lemma}[Transport]\label{lem:transport}
Let $g\in\Sn_4$, $K=g^{-1}Hg$, and let
$\chi^g(g^{-1}\sigma g)=\chi(\sigma)$.  If
$A^g_{ij}=A_{g(i),g(j)}$, then
\[
 d_{\chi^g}^{K}(A^g)=d_\chi^H(A),\qquad \per A^g=\per A,
 \qquad A\psd\Longrightarrow A^g\psd.
\]
\end{lemma}

\begin{proof}
Simultaneous permutation of rows and columns preserves positivity and the
permanent.  For a permutation monomial, changing $i$ to $g(i)$ gives
\[
 \prod_i A^g_{i,(g^{-1}\sigma g)(i)}
 =\prod_j A_{j,\sigma(j)}.
\]
Summing proves the assertion.
\end{proof}

Thus it suffices to prove dominance for the thirty-seven rows of the eleven
fixed representatives.

\section{Correlation coordinates and the routine rows}

\subsection{Four cycle aggregates}

For~\eqref{eq:corr-six} define
\begin{align}
 \mathsf T={}&|a|^2+|b|^2+|c|^2+|d|^2+|e|^2+|f|^2,\label{eq:T}\\
 \mathsf D={}&|a|^2|f|^2+|b|^2|e|^2+|c|^2|d|^2,\label{eq:D}\\
 \mathsf C={}&2\Rea(a d\bar b+a e\bar c+b f\bar c+d f\bar e),\label{eq:C}\\
 \mathsf F={}&2\Rea(a d f\bar c+a e\bar f\bar b+b\bar d e\bar c).
 \label{eq:F}
\end{align}
Enumerating the cycle types of $\Sn_4$ gives
\begin{equation}\label{eq:per-aggregate}
 \per A=1+\mathsf T+\mathsf D+\mathsf C+\mathsf F.
\end{equation}
The five normalized $S_4$ rows in the order~\eqref{eq:s4-table} are
\begin{align}
 p_{[4]}&=1+\mathsf T+\mathsf D+\mathsf C+\mathsf F,\nonumber\\
 p_{[1^4]}&=1-\mathsf T+\mathsf D+\mathsf C-\mathsf F,\nonumber\\
 p_{[31]}&=1+\tfrac13\mathsf T-\tfrac13\mathsf D-\tfrac13\mathsf F,\label{eq:s4-rows}\\
 p_{[211]}&=1-\tfrac13\mathsf T-\tfrac13\mathsf D+\tfrac13\mathsf F,\nonumber\\
 p_{[22]}&=1+\mathsf D-\tfrac12\mathsf C.\nonumber
\end{align}

\begin{lemma}[Scalar correlation inequalities]\label{lem:scalar-toolbox}
If $A$ in~\eqref{eq:corr-six} is positive semidefinite, then
\begin{gather}
 0\leq\mathsf T,\quad 0\leq\mathsf D,\quad
 \mathsf D\leq\frac{\mathsf T^2}{4},\quad
 -2\mathsf D\leq\mathsf F\leq2\mathsf D\leq\mathsf T,
 \label{eq:scalar-basic}\\
 \frac{\mathsf T^2}{2}\leq\mathsf T+\frac32\mathsf C,
 \qquad 1\leq\per A.\label{eq:moment-bound}
\end{gather}
Every $3\times3$ principal permanent is at most $\per A$.
\end{lemma}

\begin{proof}
The inequalities $|a|^2,\ldots,|f|^2\leq1$ follow from the $2\times2$
principal minors.  The first four estimates in~\eqref{eq:scalar-basic} are
then AM--GM and
$2|\Rea(xy)|\leq|x|^2+|y|^2$, paired over opposite edges.

For the moment inequality, let $\lambda_1,\dots,\lambda_4\geq0$ be the
eigenvalues of $A$.  Direct trace expansion gives
\[
 \sum\lambda_i=4,\quad
 \sum\lambda_i^2=4+2\mathsf T,\quad
 \sum\lambda_i^3=4+6\mathsf T+3\mathsf C.
\]
The exact identity
\[
 \left(\sum_i\lambda_i\right)\left(\sum_i\lambda_i^3\right)
 -\left(\sum_i\lambda_i^2\right)^2
 =\sum_{i<j}\lambda_i\lambda_j(\lambda_i-\lambda_j)^2
\]
implies~\eqref{eq:moment-bound}'s first inequality.  Combining it with
$\mathsf D\leq\mathsf T^2/4$ and $\mathsf F\geq-2\mathsf D$ gives
$1\leq1+\mathsf T+\mathsf D+\mathsf C+\mathsf F$.

Finally, the permanent version of the $3+1$ Fischer contraction can be
verified directly from~\eqref{eq:per-aggregate}: after subtracting a
principal $3\times3$ permanent, the remaining expression is the sum of the
corresponding $3\times3$ determinant inequality and squared edge terms.
Applying this to each of the four principal triples proves the last claim.
\end{proof}

We also need two symmetric $2+2$ contractions.  Put
\begin{align}
 R_+&:=1+|b|^2+|e|^2+\mathsf D
       +2\Rea(a d f\bar c),\label{eq:Rplus}\\
 R_\wedge&:=(1-|b|^2)(1-|e|^2)+|af-c\bar d|^2.\label{eq:Rwedge}
\end{align}

\begin{lemma}[Block contractions]\label{lem:block-contractions}
For a positive-semidefinite correlation matrix,
\begin{equation}\label{eq:block-bounds}
 R_+\leq\per A,\qquad R_\wedge\leq\per A,\qquad
 1+|a|^2|f|^2+|c|^2|d|^2\leq\per A.
\end{equation}
\end{lemma}

\begin{proof}
For the split $\{1,3\}\mid\{2,4\}$, take Gram vectors
$x_1,x_2,x_3,x_4$ for $A$.  Orthogonally remove the component of $x_3$
along $x_1$ and apply Cauchy--Schwarz to the resulting two pairs.  Expansion
gives the first inequality.  Applying exterior Cauchy--Schwarz to
$x_1\wedge x_3$ and $x_2\wedge x_4$ gives
\[
 |af-c\bar d|^2\leq(1-|b|^2)(1-|e|^2),
\]
and the same block expansion gives the second inequality.  Relabelling the
split twice yields
$(1+|a|^2)(1+|f|^2)\leq\per A$ and
$(1+|c|^2)(1+|d|^2)\leq\per A$; comparison of the two opposite-edge
products gives the third inequality.
\end{proof}

\subsection{Explicit dispatch of the routine character rows}

The following proposition records the finite analytic part of the proof.
Its table is included to make clear that the only rows not covered are the
two non-real $A_4$ rows.

\begin{proposition}[Thirty-five routine rows]\label{prop:routine-rows}
For every representative subgroup in Lemma~\ref{lem:subgroup-classification},
all normalized irreducible rows other than $\chi_\omega$ and
$\bar\chi_\omega$ of $A_4$ are bounded above by $\per A$.
\end{proposition}

\begin{proof}
The cyclic, Klein-four, and $S_3$ expansions are as follows.  In the $C_3$
and $S_3$ rows put $z=ad\bar b$.
\begin{center}
\begin{tabular}{@{}ll@{}}
\toprule
group & normalized row values\\
\midrule
$1$ & $1$\\
$C_2^{\rm t}$ & $1\pm|a|^2$\\
$C_2^{\rm d}$ & $1\pm|a|^2|f|^2$\\
$C_3$ & $1+2\Rea z$,\quad $1-\Rea z\pm\sqrt3\Ima z$\\
$V_4^{\rm n}$ & $1\pm r\pm s\pm t$ with the four character sign patterns,\\
& $r=|a|^2|f|^2$, $s=|b|^2|e|^2$, $t=|c|^2|d|^2$\\
$V_4^{\rm d}$ & $(1\pm|a|^2)(1\pm|f|^2)$\\
$S_3$ & $1+|a|^2+|b|^2+|d|^2+2\Rea z$,\\
& $1-|a|^2-|b|^2-|d|^2+2\Rea z$,\quad $1-\Rea z$\\
\bottomrule
\end{tabular}
\end{center}
These expressions follow by listing the elements of the corresponding
subgroup.  The $3\times3$ principal determinant inequalities, the principal
permanent contraction in Lemma~\ref{lem:scalar-toolbox}, and
$|a|^2,\ldots,|f|^2\leq1$ prove their bounds.  For example, the three $C_3$
rows are precisely the normalized one-dimensional rows of the principal
$3\times3$ correlation block, while the three $S_3$ rows are its trivial,
sign, and standard rows.

For $C_4$, put
\[
 r=|b|^2|e|^2,\qquad z=adf\bar c.
\]
Its four rows are
\begin{equation}\label{eq:c4-rows}
 1+r+2\Rea z,\quad 1-r-2\Ima z,\quad
 1+r-2\Rea z,\quad 1-r+2\Ima z.
\end{equation}
The estimates
\[
 2|\Rea z|,\ 2|\Ima z|
 \leq |a|^2|f|^2+|c|^2|d|^2
\]
together with Lemma~\ref{lem:block-contractions} prove all four bounds.

For $D_8$, write $p=|b|^2$, $q=|e|^2$.  The five rows are
\begin{align}
 R_+,&\quad
 (1+p)(1+q)-|af+c\bar d|^2,\nonumber\\
 &(1-p)(1-q)-|af-c\bar d|^2,\quad
 R_\wedge,\quad 1-pq.\label{eq:d8-rows}
\end{align}
The first and fourth are~\eqref{eq:block-bounds}; the second is no larger
than $R_+$ by direct expansion; the third and fifth are at most $1$, hence at
most $\per A$ by Lemma~\ref{lem:scalar-toolbox}.

For $S_4$, substitute~\eqref{eq:scalar-basic}--\eqref{eq:moment-bound} into
the five expressions~\eqref{eq:s4-rows}.  The sign row uses
$\mathsf T+\mathsf F\geq0$; the $[31]$ and $[22]$ rows use the moment
inequality; and the $[211]$ row is at most $1$ because
$\mathsf F\leq\mathsf T$ and $\mathsf D\geq0$.

Finally, the trivial $A_4$ row is the average of the trivial and sign $S_4$
rows, while its three-dimensional row is the average of the $[31]$ and
$[211]$ rows.  This proves the asserted thirty-five rows.
\end{proof}

\section{The exceptional non-real \texorpdfstring{$A_4$}{A4} gap}

Let
\begin{equation}\label{eq:Buvw}
 B(u,v,w)=\mat{1&u&v\\\bar u&1&w\\\bar v&\bar w&1}.
\end{equation}
Define
\begin{equation}\label{eq:c0}
 c_0(u,v,w)=|u|^2+|v|^2+|w|^2
 +2\Rea\bigl((1-\omega)uw\bar v\bigr)
\end{equation}
and the Hermitian matrix
\begin{equation}\label{eq:K}
 K(u,v,w)=
 \mat{
 1&(1-\omega^2)u+v\bar w&uw+(1-\omega)v\\
 (1-\omega)\bar u+w\bar v&1&(1-\omega^2)w+v\bar u\\
 \bar u\bar w+(1-\omega^2)\bar v&(1-\omega)\bar w+u\bar v&1
 }.
\end{equation}

\begin{lemma}[Character expansion]\label{lem:a4-expansion}
For the correlation matrix~\eqref{eq:corr-six},
\begin{equation}\label{eq:a4-gap}
 \per A-d_{\chi_\omega}^{A_4}(A)
 =c_0(a,b,d)+\mat{\bar c&\bar e&\bar f}
 K(a,b,d)\mat{c\\e\\f}.
\end{equation}
For $\bar\chi_\omega$, the same formula holds after an odd simultaneous
relabeling of the four indices.
\end{lemma}

\begin{proof}
List the identity, three double transpositions, and the two classes of four
three-cycles in $A_4$, insert their values from~\eqref{eq:a4-table}, and pair
inverse permutations.  Using $\omega^2=-1-\omega$ gives exactly
\eqref{eq:c0}--\eqref{eq:K}.  An odd permutation interchanges the two
three-cycle classes and therefore exchanges $\chi_\omega$ with its conjugate.
\end{proof}

The remaining task is consequently the following four-vector inequality.

\begin{theorem}[Four-vector inequality]\label{thm:four-vector}
Suppose
\begin{equation}\label{eq:completion}
 \mathcal A=\mat{B(u,v,w)&q\\q^*&1}\psd,\qquad q\in\CC^3.
\end{equation}
Then
\begin{equation}\label{eq:four-vector}
 c_0(u,v,w)+q^*K(u,v,w)q\geq0.
\end{equation}
\end{theorem}

The proof occupies Sections~\ref{sec:normalization}--\ref{sec:Ppositive}.
Before giving it, we record the matrix certificate that motivated the
geometric formulation.

\begin{lemma}[Ellipsoid-to-adjugate certificate]\label{lem:ellipsoid}
Let $B,K\in M_m(\CC)$ with $B\psd$ and $K=K^*$, and let $c\in\RR$.  If
\[
 \mat{B&q\\q^*&1}\psd\quad\Longrightarrow\quad c+q^*Kq\geq0
\]
for every $q\in\CC^m$, then
\begin{equation}\label{eq:adj-certificate}
 (\det B)K+c\,\adj(B)\psd.
\end{equation}
\end{lemma}

\begin{proof}
If $\det B=0$, the choice $q=0$ gives $c\geq0$ and
$\adj(B)\psd$, so~\eqref{eq:adj-certificate} follows.  If $\det B>0$, then
$B\pd$ and the block condition is $q^*B^{-1}q\leq1$.  Given $y\neq0$, set
\[
 q=\frac{y}{\sqrt{y^*B^{-1}y}}.
\]
The hypothesis, multiplied by $y^*B^{-1}y$, says
$y^*(K+cB^{-1})y\geq0$.  Hence $K+cB^{-1}\psd$.  Multiplication by
$\det B>0$ and the identity $\adj(B)=(\det B)B^{-1}$ prove the claim.
\end{proof}

Applied to~\eqref{eq:c0}--\eqref{eq:K}, this gives the geometric spectral
certificate
\begin{equation}\label{eq:a4-certificate}
 (\det B(u,v,w))K(u,v,w)+c_0(u,v,w)\adj(B(u,v,w))\psd.
\end{equation}

\section{Residual-vector normalization}\label{sec:normalization}

Assume first that $q_1q_2q_3\neq0$.  The Schur complement of the final
diagonal entry in~\eqref{eq:completion} is
$B-qq^*\psd$.  Congruence by
$\diag(\bar q_1^{-1},\bar q_2^{-1},\bar q_3^{-1})$ gives the residual Gram
matrix
\begin{equation}\label{eq:residual-gram}
 R=\mat{
 n_1&h_{12}&h_{13}\\
 \bar h_{12}&n_2&h_{23}\\
 \bar h_{13}&\bar h_{23}&n_3
 }\psd,
\end{equation}
where
\begin{equation}\label{eq:residual-coordinates}
 n_i=\frac{1-|q_i|^2}{|q_i|^2},\qquad
 h_{12}=\frac{u}{q_1\bar q_2}-1,\quad
 h_{13}=\frac{v}{q_1\bar q_3}-1,\quad
 h_{23}=\frac{w}{q_2\bar q_3}-1.
\end{equation}
Thus $R$ is the Gram matrix of three residual vectors $w_1,w_2,w_3$ with
$n_i=\norm{w_i}^2$ and $h_{ij}=\ip{w_i}{w_j}$.

We now define the normalized geometric polynomial.  Put
\begin{align}
 E={}&n_1n_2+|h_{12}|^2+n_1n_3+|h_{13}|^2+n_2n_3+|h_{23}|^2,\nonumber\\
 x_{12,13}={}&n_1h_{23}+h_{13}\bar h_{12},\nonumber\\
 x_{12,23}={}&h_{12}h_{23}+n_2h_{13},\nonumber\\
 x_{13,23}={}&n_3h_{12}+h_{13}\bar h_{23},\nonumber\\
 S={}&E+2\Rea(x_{12,13}+x_{12,23}+x_{13,23}),\label{eq:tensor-terms}\\
 S_\omega={}&E+2\Rea(\omega^2x_{12,13}+\omega x_{12,23}
                     +\omega^2x_{13,23}),\nonumber\\
 \tau={}&h_{13}\bar h_{12}\bar h_{23},\nonumber\\
 \Gamma={}&n_3|h_{12}|^2+n_2|h_{13}|^2+n_1|h_{23}|^2
              +3\Rea\tau+\sqrt3\Ima\tau.\nonumber
\end{align}
Finally set
\begin{equation}\label{eq:geometric-P}
 \mathcal P(R)=24+6\left(n_1+n_2+n_3
 +2\Rea(h_{12}+h_{13}+h_{23})\right)+2S-S_\omega+\Gamma.
\end{equation}

\begin{lemma}[Exact normalization identity]\label{lem:normalization-identity}
With~\eqref{eq:residual-coordinates},
\begin{equation}\label{eq:normalization-identity}
 c_0(u,v,w)+q^*K(u,v,w)q
 =|q_1q_2q_3|^2\mathcal P(R).
\end{equation}
\end{lemma}

\begin{proof}
Equations~\eqref{eq:residual-coordinates} give
\[
 u=q_1\bar q_2(1+h_{12}),\quad
 v=q_1\bar q_3(1+h_{13}),\quad
 w=q_2\bar q_3(1+h_{23}),
\]
and $|q_i|^2(1+n_i)=1$.  Substitute these identities in
\eqref{eq:c0} and~\eqref{eq:K}, use
$\omega^2+\omega+1=0$, and collect the coefficients of the three
normalization equations.  The diagonal quadratic terms give $E$, the three
mixed symmetric-tensor terms give $S$ and $S_\omega$, and the cubic terms
give $\Gamma$.  The result is~\eqref{eq:normalization-identity}.
\end{proof}

Thus Theorem~\ref{thm:four-vector} follows from the next result.

\begin{theorem}[Geometric polynomial]\label{thm:Ppositive}
For every Hermitian matrix $R$ of the form~\eqref{eq:residual-gram},
\[
 R\psd\quad\Longrightarrow\quad\mathcal P(R)\geq0.
\]
\end{theorem}

\section{Positivity of the geometric polynomial}\label{sec:Ppositive}

\subsection{The rank-one boundary}

Suppose first that $R$ has rank one.  Write
\[
 n_1=|x|^2,\quad n_2=|y|^2,\quad n_3=|z|^2,\qquad
 h_{12}=\bar xy,\quad h_{13}=\bar xz,\quad h_{23}=\bar yz.
\]
Define
\begin{align}
 R_0(x,y,z)={}&4+|x+y+z|^2+|xy+xz+yz|^2\nonumber\\
 &\quad-(|xy|^2+|xz|^2+|yz|^2)+|xyz|^2,\label{eq:R0}\\
 \Phi(x,y,z)={}&xy+\omega xz+\omega^2yz.
\end{align}

\begin{lemma}[Rank-one sum of squares]\label{lem:rank-one-sos}
For all $x,y,z\in\CC$,
\begin{equation}\label{eq:P-rank-one}
 \mathcal P(R)=2\bigl(3R_0(x,y,z)+|\Phi(x,y,z)|^2\bigr)\geq0.
\end{equation}
\end{lemma}

\begin{proof}
Regard $R_0$ as a quadratic in $z$.  Put
\[
 K_0=|1+\bar xy|^2,\quad
 L_0=\overline{x+y}+\overline{xy}(x+y),\quad
 C_0=4+|x+y|^2.
\]
Direct expansion gives
\[
 R_0=K_0|z|^2+2\Rea(L_0z)+C_0
\]
and the exact discriminant identity
\[
 C_0K_0-|L_0|^2=4(1+\Rea(\bar xy))^2\geq0.
\]
Completing the square proves $R_0\geq0$.  A second direct expansion using
$1+\omega+\omega^2=0$ gives the equality in~\eqref{eq:P-rank-one}; its
right-hand side is nonnegative.
\end{proof}

\subsection{The positive-definite region}

It remains to treat a positive-definite residual Gram matrix.  Write
\begin{equation}\label{eq:G-data}
 a=n_1>0,\quad b=n_2,\quad c=n_3,\quad
 h_{12}=g=r-is,\quad p=h_{13},\quad q=h_{23},
\end{equation}
and put
\begin{equation}\label{eq:Gram-G-delta}
 G=\mat{a&g\\\bar g&b},\qquad
 \delta=\det G=ab-r^2-s^2>0,\qquad h=\mat{p\\q}.
\end{equation}
There is a unique $\zeta\in\CC^2$ such that $h=G\zeta$.  Positivity of
$R$ tested on $(-\zeta,1)$ gives the Schur slack
\begin{equation}\label{eq:schur-slack}
 \zeta^*G\zeta\leq c.
\end{equation}

Set
\begin{align}
 \lambda={}&6+a+b+r^2+s^2+5r-\sqrt3s,\label{eq:lambda}\\
 P_0={}&24+6(a+b+2r)+ab+r^2+s^2,\label{eq:P0}\\
 \ell={}&\mat{
 6+(2-\omega)g+(2-\omega^2)b\\
 6+(2-\omega)a+(2-\omega^2)\bar g},\label{eq:ell}\\
 M={}&\mat{
 1+b&2-\omega+(1-\omega)g\\
 2-\omega^2+(1-\omega^2)\bar g&1+a},\label{eq:M}\\
 T={}&\lambda I_2+MG.\label{eq:Tmatrix}
\end{align}

\begin{lemma}[Quadratic decomposition]\label{lem:quadratic-decomposition}
With the preceding notation,
\begin{align}
 \mathcal P(R)
 ={}&P_0+2\Rea(\ell^*h)+\lambda c+\Rea(h^*Mh)\label{eq:P-quadratic}\\
 ={}&Q(\zeta)+\lambda(c-\zeta^*G\zeta),\nonumber
\end{align}
where
\begin{equation}\label{eq:Qzeta}
 Q(\zeta)=P_0+2\Rea((G\ell)^*\zeta)+\zeta^*(GT)\zeta.
\end{equation}
Moreover $GT$ is Hermitian.
\end{lemma}

\begin{proof}
Substitution of~\eqref{eq:G-data} in~\eqref{eq:geometric-P} gives the first
identity.  Replacing $h$ by $G\zeta$, using $G=G^*$ and
$T=\lambda I+MG$, gives the second.  Finally
$GT=\lambda G+GMG$ is Hermitian because $G$ and $M$ are Hermitian.
\end{proof}

\begin{lemma}[Positivity of the scalar part]\label{lem:lambda-positive}
Under $a,b\geq0$ and $r^2+s^2\leq ab$, one has $\lambda>0$.
\end{lemma}

\begin{proof}
Let $\rho=\sqrt{r^2+s^2}$.  Cauchy--Schwarz and AM--GM give
\[
 5r-\sqrt3s\geq-2\sqrt7\rho,\qquad a+b\geq2\rho.
\]
Consequently
\[
 \lambda\geq(\rho-(\sqrt7-1))^2+2(\sqrt7-1)>0.
\]
\end{proof}

The remaining algebra is organized by the Gram defect $\delta$.  Define
\begin{equation}\label{eq:Ndef}
 N=\Rea\left(P_0\det T-\ell^*G\adj(T)\ell\right).
\end{equation}

\begin{lemma}[Hessian and minimum certificates]\label{lem:hessian-certs}
For $a>0$ and $\delta\geq0$,
\begin{equation}\label{eq:detT-expansion}
 a^3\Rea\det T=D_0+\delta D_1+\delta^2D_2,
\end{equation}
where
\begin{align}
 D_0&=a^3\lambda_0H_0>0,\nonumber\\
 D_1&=a\bigl((4+11a)(r^2+s^2)+a(20+a)r
 -\sqrt3a(4+3a)s+a(18+10a+3a^2)\bigr)>0,\label{eq:Dcoeffs}\\
 D_2&=a(a+1)(a+2)>0.
\end{align}
Here $\lambda_0$ denotes~\eqref{eq:lambda} with
$b=(r^2+s^2)/a$, and
\[
 H_0=2\lambda_0-6+4(r^2+s^2)>0.
\]
Furthermore
\begin{equation}\label{eq:N-expansion}
 a^3N=N_0+\delta N_1+\delta^2N_2+\delta^3N_3,
\end{equation}
with $N_0\geq0$ and $N_1,N_2,N_3>0$.
\end{lemma}

\begin{proof}
The determinant identity follows by expanding~\eqref{eq:Tmatrix} after
$b=(r^2+s^2+\delta)/a$.  Positivity of $H_0$ follows from
Lemma~\ref{lem:lambda-positive} and the sharper estimate obtained from the
same two inequalities.  For $D_1$, complete squares in $(r,s)$; the residual
is
\[
 4a(72+126a+94a^2+26a^3)>0.
\]
The assertion for $D_2$ is immediate.

The expansion~\eqref{eq:N-expansion} is obtained in the same way from
\eqref{eq:Ndef}.  On the boundary $\delta=0$, the first two residual vectors
are collinear.  Lemma~\ref{lem:rank-one-sos}, viewed as a quadratic in the
third scalar coordinate, implies that its discriminant is nonpositive.  In
the notation above this gives
\[
 N_0=a^3\lambda_0(P_0H_0-|L_0'|^2)\geq0
\]
for the corresponding line coefficient $L_0'$.

The remaining coefficients are explicit.  First,
\[
 N_3=(a+2)(a^2+6)>0.
\]
Next
\begin{align*}
 N_2={}&36(r^2+s^2)
 +a(120+84r-36\sqrt3s+36(r^2+s^2))\\
 &+a^2(60+24r-24\sqrt3s+16(r^2+s^2))\\
 &+a^3(16+2(r^2+s^2))+2a^4.
\end{align*}
Completing squares gives the strictly positive residual
\[
 16a(a^6+16a^5+112a^4+318a^3+588a^2+936a+1080).
\]

For $N_1$, put $X=r$, $Y=\sqrt3s$ and
$\nu=X^2+Y^2/3$.  Then
\begin{equation}\label{eq:N1-B}
 N_1=36\nu^2+aB_1+a^2B_2+a^3B_3+a^4B_4+6a^5,
\end{equation}
where
\begin{align*}
 B_1&=\nu(240+168X-72Y+54\nu),\\
 B_2&=432+624X-240Y+492X^2+108Y^2-168XY\\
    &\qquad +(156X-84Y)\nu+38\nu^2,\\
 B_3&=168+72X-168Y+54X^2+46Y^2-84XY\\
    &\qquad +(20X-24Y)\nu+6\nu^2,\\
 B_4&=60+24X-24Y+14\nu.
\end{align*}
Appendix~\ref{app:sos} gives exact certificates for
$B_1\geq0$, $B_2>0$, $B_4>0$, and
$B_2+aB_3+a^2B_4>0$.  Hence~\eqref{eq:N1-B} is strictly positive.
\end{proof}

\begin{proposition}[Positive-definite residual Gram matrices]
If $R\pd$, then $\mathcal P(R)>0$.
\end{proposition}

\begin{proof}
Equations~\eqref{eq:detT-expansion} and~\eqref{eq:Dcoeffs} show
$\Rea\det T>0$.  The leading entry of $GT$ equals a positive rank-one
boundary value plus $\delta(a+1)$, hence it is positive.  Since $GT$ is
Hermitian and
\[
 \det(GT)=\delta\det T>0,
\]
the $2\times2$ matrix $GT$ is positive definite.

The critical point of~\eqref{eq:Qzeta} is
\[
 \zeta_0=-\frac{\adj(T)\ell}{\Rea\det T},
\]
and completing the Hermitian square gives
\[
 \min_\zeta Q(\zeta)=\frac{N}{\Rea\det T}.
\]
By~\eqref{eq:N-expansion}, $N>0$ when $\delta>0$.  Therefore
$Q(\zeta)>0$ for all $\zeta$.  Lemma~\ref{lem:lambda-positive} and the Schur
slack~\eqref{eq:schur-slack} now give
\[
 \mathcal P(R)=Q(\zeta)+\lambda(c-\zeta^*G\zeta)>0.
\]
\end{proof}

\subsection{Closure of the Gram cone}

\begin{proof}[Proof of Theorem~\ref{thm:Ppositive}]
For $R\psd$ and $t>0$, the matrix $R+tI$ is positive definite, and hence
$\mathcal P(R+tI)>0$ by the preceding proposition.  Since
$\mathcal P$ is a polynomial in the real and imaginary parts of the entries,
letting $t\downarrow0$ gives $\mathcal P(R)\geq0$.
\end{proof}

\begin{proof}[Proof of Theorem~\ref{thm:four-vector}]
If all coordinates of $q$ are nonzero, the residual matrix
\eqref{eq:residual-gram} is positive semidefinite.  Theorem~\ref{thm:Ppositive}
and~\eqref{eq:normalization-identity} give~\eqref{eq:four-vector}.

For arbitrary $q$, choose phases $p_i$ of modulus one, taking
$p_i=q_i/|q_i|$ when $q_i\neq0$ and $p_i=1$ otherwise.  Convexly interpolate
the block matrix~\eqref{eq:completion} with the rank-one correlation
completion generated by $(p_1,p_2,p_3,1)$.  Its new last column is
$q_i(t)=(1-t)q_i+tp_i$, which is nonzero for $0<t<1$.  The interpolated
matrix remains positive semidefinite, so the already proved case applies.
Letting $t\downarrow0$ and using continuity proves the result.
\end{proof}

\section{Completion of the proof}

\begin{proposition}[The two non-real $A_4$ rows]\label{prop:a4-rows}
For every $4\times4$ positive-semidefinite correlation matrix $A$,
\[
 d_{\chi_\omega}^{A_4}(A)\leq\per A,\qquad
 d_{\bar\chi_\omega}^{A_4}(A)\leq\per A.
\]
\end{proposition}

\begin{proof}
Write $A$ in the block form~\eqref{eq:completion} with
$B=B(a,b,d)$ and $q=(c,e,f)^T$.  Lemma~\ref{lem:a4-expansion} and
Theorem~\ref{thm:four-vector} give the first inequality.  An odd simultaneous
relabeling preserves positivity and the permanent and interchanges the two
non-real characters, giving the second.
\end{proof}

\begin{proof}[Proof of Theorem~\ref{thm:main}]
By Lemma~\ref{lem:correlation-reduction}, it is enough to consider
correlation matrices.  Lemma~\ref{lem:subgroup-classification} and the
transport Lemma~\ref{lem:transport} reduce an arbitrary subgroup and
irreducible representation to one of the thirty-seven registered rows.
Proposition~\ref{prop:routine-rows} proves thirty-five of them and
Proposition~\ref{prop:a4-rows} proves the remaining two.  Lemma~\ref{lem:reality}
shows that the compared quantities are real.  Finally,
Lemma~\ref{lem:scaling} returns from correlation matrices to the full
positive-semidefinite cone.
\end{proof}

\appendix

\section{Exact scalar and Gram certificates}\label{app:sos}

This appendix records the non-obvious polynomial certificates used in
Lemma~\ref{lem:hessian-certs}.  They contain no numerical approximation.
Recall $\nu=X^2+Y^2/3$.

\subsection{Direct square completions}

The coefficient $B_2$ has the exact decomposition
\begin{align*}
 B_2={}&38\left(\nu+\frac{39}{19}X-\frac{21}{19}Y+\frac32\right)^2
 +\frac{448}{19}\left(Y+\frac3{32}X-\frac{1083}{448}\right)^2\\
 &+\frac{3483}{16}\left(X+\frac{2137}{2322}\right)^2
 +\frac{131797}{5418}>0.
\end{align*}
Likewise
\begin{align*}
 B_3+252={}&6\left(\nu+\frac53X-2Y-\frac12\right)^2
 +24\left(Y-\frac{11}{12}X-\frac{15}{4}\right)^2\\
 &+\frac{139}{6}\left(X-\frac{249}{139}\right)^2
 +\frac{1851}{278}>0.
\end{align*}
Further elementary completions give
\[
 B_1\geq\frac{112}{3}\nu\geq0,\qquad
 B_4\geq\frac{132}{7}>0.
\]

\subsection{The conditional middle-coefficient certificate}

Set
\[
 \Sigma=4B_2B_4+252B_3.
\]
An exact Gram decomposition is
\begin{equation}\label{eq:sigma-gram}
 \Sigma=2128\nu\left(\nu+\frac{387}{133}X-
 \frac{261}{133}Y+\frac{77}{20}\right)^2
 +\frac1{59850}m^TQm,
\end{equation}
where $m=(1,Y,Y^2,X,XY,X^2)^T$ and
\begingroup
\scriptsize
\[
Q=\mat{
8739057600&-4231634400&478800000&6265576800&-1795500000&47880000\\
-4231634400&2519799114&-385003080&-2526627600&807975000&239400000\\
478800000&-385003080&79782560&281630160&-81252000&-35910000\\
6265576800&-2526627600&281630160&10931106942&-2359670040&2303554680\\
-1795500000&807975000&-81252000&-2359670040&685371360&-243756000\\
47880000&239400000&-35910000&2303554680&-243756000&1122611040
}.
\]
\endgroup
Exact rational $LDL^T$ elimination has positive pivots
\begin{align*}
&8739057600,\quad \frac{79557301866}{169},\quad
\frac{1112696849800960}{299087601},\\
&\frac{2624099206828303859949}{869294413907},\quad
\frac{2531308970961028469920080}{115380521779374043},\\
&\frac{1406694633431943161674826341030}
{10547120712337618624667}.
\end{align*}
Thus $Q\pd$ and $m^TQm>0$ because the first coordinate of $m$ is one.
Equation~\eqref{eq:sigma-gram} proves $\Sigma>0$.

If $B_3\geq0$, then $B_2+aB_3+a^2B_4>0$ is immediate.  If $B_3<0$, the
identity
\[
 4B_2B_4-B_3^2=(4B_2B_4+252B_3)+(-B_3)(B_3+252)
\]
shows that the discriminant of the quadratic
$B_2+aB_3+a^2B_4$ is negative.  Hence that quadratic is positive for every
$a>0$.  This completes the proof that $N_1>0$.

\section{Character-row audit}

For convenience, the logical dispatch is summarized below.
\begin{longtable}{@{}p{0.20\textwidth}p{0.25\textwidth}p{0.47\textwidth}@{}}
\toprule
subgroup types & rows & analytic certificate\\
\midrule
\endfirsthead
\toprule
subgroup types & rows & analytic certificate\\
\midrule
\endhead
$1,C_2,C_3,V_4,S_3$ & 19 rows & principal $2\times2$ and $3\times3$
minor inequalities and principal-permanent contraction\\
$C_4$ & 4 rows & opposite-edge AM--GM plus the two block contractions\\
$D_8$ & 5 rows & symmetric and exterior $2+2$ block contractions\\
$S_4$ & 5 rows & eigenvalue moment identity, cycle bounds, and
$\per A\geq1$\\
$A_4$ & trivial and $3$-dimensional & averages of already proved $S_4$ rows\\
$A_4$ & $\chi_\omega,\bar\chi_\omega$ & residual-vector geometric
polynomial and Theorem~\ref{thm:four-vector}\\
\bottomrule
\end{longtable}

The row count in the first entry is
$1+2+2+3+4+4+3=19$; adding $4+5+5+4$ gives $37$.  The $A_4$ line contains
four rows in total, two routine and two geometric.

\section{Correspondence with the Lean formalization}

The proof has been formalized in Lean 4.19.0 with Mathlib 4.19.0.  The main
mathematical correspondences are:
\begin{center}
\begin{tabular}{@{}p{0.40\textwidth}p{0.52\textwidth}@{}}
\toprule
manuscript component & Lean modules\\
\midrule
definitions and normalization & \texttt{Basic}, \texttt{DiagonalScaling},
\texttt{CorrelationLift}\\
subgroups and character completeness & \texttt{SubgroupRegistry},
\texttt{CharacterTables}, realization and completeness registries\\
routine analytic rows & \texttt{LowOrderRegistry}, \texttt{C4TableBridge},
\texttt{D8TableBridge}, \texttt{S4TableBridge}\\
$A_4$ expansion and certificate & \texttt{A4Certificate}, \texttt{A4Gap},
\texttt{A4GeometricBridge}\\
residual normalization & \texttt{A4GeometricNormalization},
\texttt{A4GeometricFourVector}\\
rank-one boundary & \texttt{A4GeometricRankOne},
\texttt{A4GeometricCollinear}\\
Hessian and minimum & \texttt{A4GeometricHessian},
\texttt{A4GeometricQuadratic}, \texttt{A4GeometricMinimum}\\
exact polynomial certificates & \texttt{A4GeometricGeneral},
\texttt{A4GeometricGram}, \texttt{A4GeometricAssembly}\\
PSD closure and final theorem & \texttt{A4GeometricPSD},
\texttt{FinalAssembly}, \texttt{NormalizedPermanentalDominance}\\
\bottomrule
\end{tabular}
\end{center}

The public formal endpoints are
\begin{verbatim}
PermanentalDominance.N4.normalized_permanental_dominance_n4
PermanentalDominance.N4.normalized_permanental_dominance_n4_real
\end{verbatim}
and the central geometric endpoints are
\begin{verbatim}
A4GeometricPSD.geometricP_nonneg_of_posSemidef
A4GeometricFourVector.fourVectorProperty
A4GeometricBridge.certificate_posSemidef_geometric
\end{verbatim}

\end{document}
